\begin{table}[htbp]\centering
\caption{Efeitos fixos de condado vs.\ painel de duas vias (condado $+$ ano)}\label{tab:q3_duas_vias}
\small
\begin{tabular}{lcc}
\toprule
 & \textbf{EF condado} & \textbf{EF duas vias} \\
 & (1) & (2) \\
\midrule
ln(prob.\ pris\~ao) $\beta_1$ & -0,395*** & -0,360*** \\
 & (0,034) & (0,033) \\
ln(prob.\ condena\c{c}\~ao) & -0,309*** & -0,286*** \\
 & (0,022) & (0,021) \\
ln(prob.\ encarc.) & -0,211*** & -0,184*** \\
 & (0,033) & (0,033) \\
ln(senten\c{c}a m\'edia) & 0,032 & -0,003 \\
 & (0,026) & (0,027) \\
ln(pol\'icia p.c.) $\beta_2$ & 0,422*** & 0,422*** \\
 & (0,028) & (0,027) \\
\midrule
Efeitos fixos de condado & Sim & Sim \\
Efeitos fixos de ano & N\~ao & Sim \\
Erros-padrao & Convencional & Convencional \\
N (obs.) & 630 & 630 \\
$R^2$ & 0,373 & 0,370 \\
\bottomrule
\end{tabular}
\par\vspace{2pt}{\footnotesize\itshape Erros-padrao entre parenteses. *** p$<$0,01; ** p$<$0,05; * p$<$0,10. Dep.: ln(taxa de crime). $R^2$ \textit{within}. Os efeitos fixos de ano captam choques macro comuns a todos os condados em cada ano.}
\end{table}
